% Target: rmp-ii-inverse-renormalization
% Formula: One MPS site splits into two sites, and one PEPS site splits into a 2x2
%   patch.
% Ink: Kernel flat and plane frames; both splits share one metric.
% Boundary: Both arrows preserve the open outer virtual and physical boundary
%   signatures.
% Source: paper TN-Review-main.tex line 1056; author source ImagesReview Section
%   II/ImagesReview Section II.tex lines 801-831
% Capabilities: kernel, plane-frame
\tenkzkernel
\[
  (a)\quad
  % one MPS site: 3-leg dot
  \begin{tenkz}[rows={wire}, cols=1, physical=up, west=open, east=open]
    \tn{}
  \end{tenkz}
  \;\longrightarrow\;
  % the split: two sites of the same bond dimension
  \begin{tenkz}[rows={wire}, cols=2, physical=up, west=open, east=open]
    \tn{} & \tn{}
  \end{tenkz}
  \qquad
  (b)\quad
  % one PEPS site: 5-leg tensor on the projected sheet
  \begin{tenkz}[lattice={1x1}, frame=plane, physical=up,
                west=open, east=open, north=open, south=open]
    \tn{}
  \end{tenkz}
  \;\longrightarrow\;
  % the would-be sqrt(D) fine-graining: a 2x2 patch
  \begin{tenkz}[lattice={2x2}, frame=plane, physical=up,
                west=open, east=open, north=open, south=open]
    \tn{} & \tn{} \\
    \tn{} & \tn{}
  \end{tenkz}
\]
